\documentclass[compress,12pt]{beamer}

\usetheme{Arguelles}
\usepackage{graphicx}
\usepackage{caption}
\usepackage[spanish,es-noshorthands]{babel}
%\usepackage{babel} 
\usepackage[pages=some]{background}
\usepackage{tikz}
\usepackage{tikz-cd}
\usepackage{amsmath,amssymb,latexsym,amscd} 
\usepackage[all,cmtip]{xy}
\usepackage{fancyhdr}
\usepackage{mathalfa}
\usepackage{mathrsfs}
\usetikzlibrary{babel}
\usepackage{hyperref}
\usepackage{ragged2e}
\usepackage{wasysym}
\usepackage{tikz}
\usetikzlibrary{arrows.meta,calc,positioning,overlay-beamer-styles,shadows.blur}
%\hypersetup{colorlinks=true,linkcolor=blue,citecolor=brown,linktocpage=true,pagebackref=true,hyperindex=true}
\pagenumbering{arabic}

\DeclareMathOperator{\op}{op}
\DeclareMathOperator{\pt}{pt}
\DeclareMathOperator{\spec}{spec}
\DeclareMathOperator{\Fit}{Fit}
\DeclareMathOperator{\Pth}{P}
\DeclareMathOperator{\Frm}{Frm}
\DeclareMathOperator{\Top}{Top}
\DeclareMathOperator{\Ord}{Ord}
\DeclareMathOperator{\Obj}{Obj}
\DeclareMathOperator{\Hom}{Hom}

\title{La V-modificación de un marco}
%\subtitle{y algunos axiomas de separación en la topología sin puntos}
\event{Seminario de Álgebra, CUCEI}
\date{\today}
\author{Juan Carlos Monter Cortés}
\institute{Universidad de Guadalajara}
\email{juan.monter2902@alumnos.udg.mx}

%\homepage{www.mywebsite.com}
\github{JCmonter}

\begin{document}

%\frame[plain]{\titlepage}

\begin{frame}[plain]
\titlepage

\vfill
\begin{center}
    \textcolor{white}{\scriptsize Investigación apoyada por SECIHTI-PROYECTO CBF2023-2024-2630}
\end{center}

\end{frame}

\begin{frame}{Contenido}
\tableofcontents
\end{frame}

\section{La construcción espacial}
\begin{frame}{Aspectos topológicos}
\begin{itemize}
    \item Sea $S\in \Top$ y $p, q\in S$, entonces 
    \[
    q \sqsubseteq p \iff q\in \overline{\{p\}}
    \]
    es el orden de especialización en $S$.
    \item Las familias de subconjuntos de $S$
    \[
    \mathcal{O}S, \quad \mathcal{C}S, \quad \mathcal{U}S, \quad \mathcal{Q}S
    \]
    abiertos, cerrados, saturados y compactos saturados.

    \item Para $E\subseteq S$, entonces 
    \[
        E^\circ, \quad \overline{E}, \quad \uparrow E
    \]
    denotan interior, cerradura y saturación de $E$.
    \end{itemize}
\end{frame}

\begin{frame}{La V-modificación de un espacio}

\end{frame}


\section{La construcción en marcos}

\begin{frame}{Cuestiones de marcos}

\end{frame}

\begin{frame}{La V-modificación de un marco}

\end{frame}

\section{¿Para qué queremos la V-modificación?}
\begin{frame}{La relación entre las V-modificaciones}

\end{frame}



\section*{\textsc{Referencias}}
\begin{frame}[allowframebreaks]
\frametitle{Bibliografía}
\begin{thebibliography}{20}\markboth{Bibliografía}{Bibliografía}

\bibitem{P.T.} P. T. Johnstone, \textit{Stone spaces}, Cambridge Studies in Advanced Mathematics, vol. 3, Cambridge University Press, Cambridge, 1982. MR 698074

%\bibitem{J.M.} J. Monter; A. Zaldívar, \textit{El enfoque locálico de las reflexiones booleanas: un análisis en la categoría de marcos} [tesis de maestría], 2022. Universidad de Guadalajara.

%\bibitem{P.S.} J. Paseka and B. Smarda, \textit{$ T_2 $-frames and almost compact frames.} Czechoslovak Mathematical Journal (1992), 42(3), 385-402.

\bibitem{J.P.} J. Picado and A. Pultr, \textit{Frames and locales: Topology without points}, Frontiers in Mathematics, Springer Basel, 2012.

\bibitem{J.P.2} J. Picado and A. Pultr, \textit{Separation in point-free topology}, Springer, 2021.

\bibitem{R.S.} RA Sexton, \textit{A point free and point-sensitive analysis of the patch assembly}, The University of Manchester (United Kingdom), 2003.

\bibitem{R.S.2} RA Sexton, \textit{Frame theoretic assembly as a unifying construct}, The University of Manchester (United Kingdom), 2000.

\bibitem{R.S.3} RA Sexton and H. Simmons, \textit{Point-sensitive and point-free patch constructions}, Journal of Pure and Applied Algebra \textbf{207} (2006), no. 2, 433-468.

\bibitem{H.S.} H. Simmons, \textit{An Introduction to Frame Theory}, lecture notes, University of Manchester. Disponible en línea en \url{https://web.archive.org/web/20190714073511/http://staff.cs.manchester.ac.uk/~hsimmons}.

\bibitem{H.S.R} H. Simmons, \textit{Regularity, fitness, and the block structure of frames.} Applied Categorical Structures 14 (2006): 1-34.

\bibitem{H.S.4} H. Simmons, \textit{The lattice theoretic part of topological separation properties}, Proceedings of the Edinburgh Mathematical Society, vol.~21, pp.~41--48, 1978.

\bibitem{H.S.V} H. Simmons, \textit{The Vietoris modifications of a frame}. Unpublished manuscript (2004), 79pp., available online at http://www. cs. man. ac. uk/hsimmons.

%\bibitem{A.W.} A. Wilansky, \textit{Between T1 and T2}, MONTHLY (1967): 261-266.

\bibitem{A.Z.} A. Zaldívar, \textit{Introducción a la teoría de marcos} [notas curso], 2025. Universidad de Guadalajara.
\end{thebibliography}
\end{frame}


\end{document}